\section{Introduccion}
\label{sec:introduccion}

Los residuos organicos representan una fraccion mayoritaria de los desechos municipales y agroindustriales, y su disposicion en vertederos es una fuente importante de emisiones de gases de efecto invernadero (GEI). La bioconversion con larvas de la mosca soldado negra (\textit{Hermetia illucens}) se ha consolidado como una alternativa de valorizacion: las larvas transforman residuos en biomasa rica en proteina y lipidos, con alta reduccion de masa y humedad del sustrato \parencite{mertenatBlackSoldierFly2019,nayakHermetiaIllucensDiptera2023}. Sin embargo, el proceso no esta exento de emisiones: la respiracion larvaria y la degradacion microbiana del lecho liberan CO$_2$, y bajo ciertas condiciones del sustrato se han reportado emisiones de CH$_4$ y N$_2$O \parencite{ermolaevGreenhouseGasEmissions2019,parodiBioconversionEfficienciesGreenhouse2020a,boakye-yiadomGreenhouseGasEmissions2022,lindbergProcessEfficiencyGreenhouse2022a}.

La cuantificacion de estas emisiones se ha hecho tipicamente con camaras estaticas o sistemas de flujo, siguiendo metodos establecidos \parencite{hutchinsonMethodsSoilAnalysis1992}, y muestra una fuerte dependencia del sustrato y las condiciones del proceso \parencite{parodiBioconversionEfficienciesGreenhouse2020a,boakye-yiadomGreenhouseGasEmissions2022}. Mas recientemente, los sensores de bajo costo han permitido el monitoreo continuo de CO$_2$ y CH$_4$ en estos sistemas \parencite{mitchellCalibrationLowCostMethane2024,kiplimoAddressingLowCostMethane2024}. No obstante, las mediciones reportan emisiones totales del sistema larvas--lecho, sin separar la contribucion de cada fuente, lo que limita su utilidad para el diseno de dietas y la operacion del proceso.

En paralelo, la modelacion matematica de la bioconversion ha avanzado hacia modelos dinamicos mecanicistas. Sobre la base de la teoria de presupuestos dinamicos de energia (DEB) \parencite{kooijmanDynamicEnergyBudget2009}, \textcite{eriksenDynamicModellingFeed2022} propuso un modelo de compartimentos que describe la asimilacion del alimento, el crecimiento estructural, la acumulacion de lipidos y la produccion de CO$_2$ de larvas de \textit{H. illucens}, luego aplicado al desempeno metabolico bajo restriccion alimentaria \parencite{eriksenMetabolicPerformanceFeed2024}. Otros trabajos han modelado el crecimiento larvario \parencite{padmanabhaComprehensiveDynamicGrowth2020}, el desempeno por dieta \parencite{laganaroGrowthMetabolicPerformance2021a}, el efecto de la temperatura \parencite{schonEffectTemperatureGrowth2024} y la estimacion de las dinamicas de GEI \parencite{rossiEstimatingDynamicsGreenhouse2024a}. A pesar de estos avances, persisten tres vacios: (i) los modelos no se han validado contra mediciones de gas que separen la fuente larvaria de la microbiana; (ii) no representan la transicion a prepupa, en la que la larva deja de alimentarse y las emisiones colapsan, aunque esta domina el final del ciclo productivo; y (iii) al omitir la degradacion microbiana del lecho, no pueden reproducir las emisiones totales medidas en un reactor por lotes.

El objetivo de este artículo es modelar y validar experimentalmente las emisiones de CO$_2$ en la bioconversion estatica con \textit{H. illucens}, distinguiendo lo que generan las larvas de lo que genera la degradacion microbiana del lecho, con un modelo construido como extension del modelo DEB de \textcite{eriksenDynamicModellingFeed2022,eriksenMetabolicPerformanceFeed2024} con tres componentes nuevos: un interruptor de prepupa que cierra la asimilacion, un modulo microbiano del lecho con un estado de biomasa cuya degradacion es de primer orden en la biomasa, limitada por el sustrato y modulada por humedad y temperatura, y el balance de gas de la camara estatica que conecta las fuentes con la medicion. La validacion explota la descomposicion de fuentes: la componente larvaria se contrasta con la tasa neta (tratamiento menos control de alimento solo) y la microbiana con los controles, en un experimento por lotes sellado con dos dietas (alimento estandar y residuos agroindustriales de frutas fermentados), cierres por triplicado (replicas tecnicas) y sensores de CO$_2$ y CH$_4$; las lecturas de CH$_4$ se usan solo como indicador cualitativo por la falta de calibracion del sensor. El articulo se organiza asi: la seccion 2 describe el sistema experimental, el modelo extendido, su solucion numerica y el analisis de incertidumbre; la seccion 3 presenta la calibracion y la validacion; la seccion 4 discute los alcances y limitaciones; y la seccion 5 concluye.
